\section{The operator abstraction}\label{sec:operators}

Nothing in Sections~\ref{sec:reduction}--\ref{sec:nonneg} requires $A$ as an explicit
matrix. The active-set loop (Algorithm~\ref{alg:elim}) together with its matrix-free CG
inner solve touches $A$ through only two operations:
\begin{itemize}
  \item a \emph{free-block application} $\mathrm{apply}_{\mathcal{F}}\colon v \mapsto
  A_{\mathcal{F}}\,v$ on the current free set, which is all CG needs to
  solve~\eqref{eq:innerspd} (and, for the equality-augmented problem, the two right-hand
  sides of Section~\ref{sec:reduction}); and
  \item a \emph{cross product} $\mathrm{cross}\colon x_{\mathcal{F}} \mapsto
  A_{\mathcal{B},\mathcal{F}}\,x_{\mathcal{F}}$ from the free block to the bound block, which
  supplies the reduced gradient $s_i = (Ax)_i - b_i$ at the bound indices $i \in \mathcal{B}$
  used by the dual-feasibility test (line~5 of Algorithm~\ref{alg:elim}), since
  $x_{\mathcal{B}} = 0$.
\end{itemize}
A full product $\mathrm{matvec}\colon x \mapsto Ax$ is convenient for a residual check but is
not otherwise required. Capturing exactly this contract as an operator protocol -- the same
device by which the critical-line algorithm reaches its covariance through a small
interface~\cite{schmelzer2026cla,cvxcla} -- lets the solver run unchanged on any object that
implements it. Several backends do, each realising the two operations at a different cost.

\paragraph{Dense.}
An explicit SPD $A \in \mathbb{R}^{n\times n}$ implements $\mathrm{apply}_{\mathcal{F}}$ by
slicing $A_{\mathcal{F}}$ and $\mathrm{cross}$ by the off-diagonal block product. This is the
reference backend: matrix-free CG then costs $O(k^2)$ per inner iteration ($k = |\mathcal{F}|$),
matching a direct Cholesky solve up to the iteration count, and is worthwhile only when the
spectrum makes CG cheap relative to a factorisation. It is the natural choice when $A$ is
already assembled and fits in memory.

\paragraph{Gram (data matrix).}
When $A = M^\top M$ for $M \in \mathbb{R}^{m\times n}$ -- the least-squares/NNLS setting -- the
backend implements the protocol straight from $M$, never forming the $n\times n$ Gram matrix:
\[
  \mathrm{apply}_{\mathcal{F}}(v) = M_{\mathcal{F}}^\top\!\left(M_{\mathcal{F}}\,v\right),
  \qquad
  \mathrm{cross}(x_{\mathcal{F}}) = M_{\mathcal{B}}^\top\!\left(M_{\mathcal{F}}\,x_{\mathcal{F}}\right),
\]
each a pair of thin products, at $O(km)$ and $O(nm)$ respectively and $O(nm)$ memory instead of
$O(n^2)$. The rank caveat is the familiar one: $M^\top M$ has rank at most $m$, so for $m < n$
the free block becomes singular once $\mathcal{F}$ outgrows the data rank -- exactly the
degeneracy the regularisation of Section~\ref{sec:conditioning} removes.

\paragraph{Factor (diagonal-plus-low-rank).}
A structured operator
\[
  A = \operatorname{diag}(d) + U\Delta U^\top,
  \qquad d \in \mathbb{R}^n_{>0},\; U \in \mathbb{R}^{n\times r},\; \Delta \in \mathbb{R}^{r\times r},
\]
implements $\mathrm{apply}_{\mathcal{F}}$ and $\mathrm{cross}$ through the diagonal and the
low-rank factor separately, and inverts the free block by the Woodbury identity
\[
  A_{\mathcal{F}}^{-1} = D_{\mathcal{F}}^{-1}
    - D_{\mathcal{F}}^{-1} U_{\mathcal{F}}\, W^{-1}\, U_{\mathcal{F}}^\top D_{\mathcal{F}}^{-1},
  \qquad W = \Delta^{-1} + U_{\mathcal{F}}^\top D_{\mathcal{F}}^{-1} U_{\mathcal{F}} \in \mathbb{R}^{r\times r},
\]
at $O(k r^2 + r^3)$ per free-block solve rather than $O(k^3)$, and $O(nr)$ memory. Here the
inner solve is \emph{direct} (the Woodbury inverse is exact), so CG is bypassed and the outer
loop consumes $A_{\mathcal{F}}^{-1}$ as a black box; the operator protocol is what lets the same
active-set loop accept either a matrix-free CG solve or a Woodbury direct solve without
modification. This backend is the subject of the second companion paper~\cite{schmelzer2026rmt}.

\paragraph{Regularised.}
The ridge split $A_\alpha = (1-\alpha)A + \alpha R^\top R$ of Section~\ref{sec:conditioning} is
itself a backend: it wraps any of the above and adds the target term to each operation. When
$A = M^\top M$ it is the Gram backend of the stacked operator
$\tilde M = \big(\sqrt{1-\alpha}\,M;\ \sqrt{\alpha}\,R\big)$, so no separate implementation is
needed. Because $A_\alpha \succ 0$ strictly, this backend also restores the $P$-matrix property
(Section~\ref{sec:problem}) when the underlying $A$ is only semidefinite.

Because Algorithm~\ref{alg:elim} and the CG inner solver reach $A$ only through
$\mathrm{apply}_{\mathcal{F}}$ and $\mathrm{cross}$, switching backend -- dense to Gram to factor
to regularised -- requires no change to the loop, the termination proof
(Theorem~\ref{thm:termination}), or the convergence analysis (Proposition~\ref{prop:cg}). The
companion papers~\cite{schmelzer2026matrixfree,schmelzer2026rmt} are, in this light, two
backends of one solver.
